% ===================================================================
%  TABLEAU N° 16 — Le grand magasin. --- Le bazar. Les jouets.
%
%  Traduction, faite en 2026, en francais standard du Quebec, sur
%  l'ido de texto/io. Regles de la colonne : voir 10-tableau-01.tex.
%  Dernier tableau du livret.
%
%  LES CINQ CHOIX PROPRES A CE TABLEAU :
%    (5)  factionnaire ........ sentinelle
%    (69) voiture d'enfant .... poussette
%         étrennes ............ cadeaux du jour de l'An
%         brancardiers ........ ambulanciers
%         bougies ............. chandelles
%
%  DEUX D'ENTRE EUX NE FONT QUE SUIVRE UN CHOIX DEJA FAIT. Les
%  « kamilisti » qui emportent les blesses du champ de bataille
%  portent la civiere du tableau 11 : ce sont donc des ambulanciers.
%  Et les « bujii » de l'arbre de Noel sont les chandelles du tableau
%  6, ou le Quebec garde le mot que la France a laisse vieillir. Une
%  colonne se tient par ses reprises autant que par ses trouvailles.
%
%  « ETRENNES » EST LE MOT QUI MANQUE ICI, ET C'EST LA FETE QUI A
%  CHANGE. Le livret de 1926 place son grand magasin a la veille du
%  jour de l'An, parce que c'etait alors le jour des cadeaux ; le mot
%  « etrennes » nommait ces cadeaux-la et ne s'emploie plus au Quebec.
%  Mais la DATE, elle, ne bouge pas : l'ido ecrit « nov yaro », le
%  sapin du renvoi 81 est bien un arbre de Noel, et les deux fetes se
%  suivent dans le texte comme elles se suivaient dans l'annee. On
%  traduit donc le mot, non la fete : cadeaux du jour de l'An.
% ===================================================================

%%K t16-apar-1 apar
\VUsaut{4.0mm}
\VUtitre{15.0pt}{60}{TABLEAU N\textsuperscript{o} 16}
\VUsaut{2.0mm}
\VUpk{13.2pt}{40}{Le grand magasin. --- Le bazar.}
\VUpk{13.2pt}{40}{Les jouets.}
\VUsaut{3.4mm}

%%K t16-01-1 p
1. --- Le nouvel an approche. Les grands magasins ont préparé leurs étalages de
\VUgras{jouets} \textsuperscript{(1)} et de \VUgras{cadeaux du jour de l'An}.

%%K t16-01-2 p
J'ai demandé à mon grand-père de m'emmener au bazar pour voir cette belle
exposition, et il a accepté.

%%K t16-02-1 p
2. --- Nous nous sommes d'abord arrêtés devant de beaux panneaux d'armes, qui
portent un \VUgras{fusil}, un \VUgras{képi}, un \VUgras{sabre}, un
\VUgras{ceinturon} et une paire d'\VUgras{épaulettes}, ou bien un
\VUgras{casque}, une \VUgras{cuirasse} \textsuperscript{(3)} et un
\VUgras{pistolet} \textsuperscript{(4)}. Tout cela m'a beaucoup plus intéressé que
les \VUgras{projecteurs} \textsuperscript{(2)}.

%%K t16-tit-1 sub
\VUsaut{3.0mm}
\VUpk{10.2pt}{40}{De belles choses à voir.}
\VUsaut{2.6mm}

%%K t16-03-1 p
3. --- Dans le même rayon, il y avait une \VUgras{sentinelle} \textsuperscript{(5)}
dans sa \VUgras{guérite}; elle semblait monter la garde devant toute une
ménagerie de carton. Que d'animaux il y avait là : un \VUgras{perroquet}
\textsuperscript{(6)} sur son perchoir, un \VUgras{rhinocéros}
\textsuperscript{(7)} avec sa corne sur le nez, un \VUgras{renne}
\textsuperscript{(8)}, une \VUgras{autruche} \textsuperscript{(9)}, un
\VUgras{singe} \textsuperscript{(10)} croquant une noix, un \VUgras{renard}
\textsuperscript{(11)} emportant une poule, un \VUgras{crocodile}
\textsuperscript{(12)} ouvrant d'énormes mâchoires, un \VUgras{chameau}
\textsuperscript{(13)}, un élégant \VUgras{lévrier} \textsuperscript{(14)}, une
\VUgras{panthère} \textsuperscript{(15)}, puis deux bêtes que je ne connaissais
pas et qui sont, m'a-t-on dit, une \VUgras{belette} \textsuperscript{(16)} et une
\VUgras{hyène} \textsuperscript{(17)}. Il y avait aussi un \VUgras{léopard}
\textsuperscript{(18)} et un \VUgras{loup} \textsuperscript{(21)}; tout près, de
mignons petits \VUgras{rats} \textsuperscript{(19)} et de minuscules
\VUgras{souris} \textsuperscript{(20)} de caoutchouc. Enfin, un
\VUgras{hippopotame} \textsuperscript{(22)} et un \VUgras{dromadaire}
\textsuperscript{(23)} bossu terminaient la collection. Je serais resté là
plusieurs heures s'il n'y avait pas eu, à côté, un splendide camp plein de
\VUgras{soldats de plomb} \textsuperscript{(29)} qui a attiré mon attention.

%%K t16-tit-2 sub
\VUsaut{4.0mm}
\VUpk{10.2pt}{40}{Soldats et guerre.}
\VUsaut{2.8mm}

%%K t16-04-1 p
4. --- Mon grand-père, qui est \VUgras{invalide} \textsuperscript{(24)} et qui a
dû quitter l'armée à cause d'une \VUgras{blessure} qui lui a valu la
\VUgras{médaille militaire}, aime beaucoup me raconter les \VUgras{campagnes}
auxquelles il a pris part. Il était heureux de m'expliquer les différents
mouvements de la petite armée miniature. Un groupe de vrais soldats qui
visitaient le bazar — un \VUgras{soldat du génie} \textsuperscript{(25)}, un
\VUgras{chasseur alpin} \textsuperscript{(26)} coiffé de son \VUgras{béret}, un
\VUgras{sergent-major} \textsuperscript{(27)} d'infanterie et un \VUgras{sergent
d'artillerie} \textsuperscript{(28)} avec son \VUgras{galon} d'or sur la manche
— s'est arrêté près de nous pour écouter ses explications.

%%K t16-05-1 p
5. --- Mon grand-père, tout fier d'avoir un auditoire aussi brillant, a
improvisé un plan de bataille complet.

%%K t16-05-2 p
La ville fortifiée, avec ses \VUgras{remparts} et ses \VUgras{fossés}, était
assiégée par l'\VUgras{armée ennemie}, qui, après un long
\VUgras{bombardement}, était prête à donner l'\VUgras{assaut}. Mais les
\VUgras{assiégés}, poussés par la faim, avaient tenté une \VUgras{sortie} vers
le \VUgras{camp} des \VUgras{assiégeants}. Après avoir repoussé l'\VUgras{attaque}
par un \VUgras{combat} acharné autour du \VUgras{campement}, les assiégeants
\VUgras{vainqueurs} avaient lancé leurs \VUgras{escadrons} à la poursuite des
assaillants \VUgras{vaincus}. Ceux-ci \VUgras{battaient en retraite} en couvrant
le \VUgras{champ de bataille} de \VUgras{morts} et de \VUgras{blessés}. Des
\VUgras{ambulanciers} emportaient ces derniers à l'\VUgras{ambulance} pour les
faire soigner par le \VUgras{chirurgien}.

%%K t16-05-3 p
Malgré la résistance héroïque de quelques \VUgras{bataillons} qui avaient formé
le \VUgras{carré}, la \VUgras{défaite} était complète, le reste de l'armée
\VUgras{fuyait} en laissant beaucoup de \VUgras{prisonniers}. L'ennemi allait
achever sa \VUgras{victoire} en occupant la ville. Ce serait peut-être la fin de
la campagne et, après un \VUgras{armistice}, il serait possible de signer un
\VUgras{traité de paix}.

%%K t16-tit-3 sub
\VUsaut{4.4mm}
\VUpk{10.2pt}{40}{Les cadeaux du jour de l'An.}
\VUsaut{2.8mm}

%%K t16-06-1 p
6. --- Les jeunes soldats, qui riaient en écoutant le récit pittoresque du
vétéran, lui ont demandé quelques autres explications. Moi, j'ai fait le tour du
rayon pour regarder une belle \VUgras{arbalète} \textsuperscript{(32)} et un
\VUgras{arc} \textsuperscript{(33)} avec sa \VUgras{flèche}
\textsuperscript{(31)}. J'y ai trouvé une autre collection d'animaux : deux
\VUgras{oiseaux chanteurs} \textsuperscript{(34)}, un \VUgras{rossignol}
mécanique et une \VUgras{fauvette} chantant à pleine gorge, et une série de
bêtes étranges : une \VUgras{chauve-souris} \textsuperscript{(35)}, un
\VUgras{boa} \textsuperscript{(36)}, une \VUgras{tortue} \textsuperscript{(37)},
un \VUgras{phoque} \textsuperscript{(38)}, une \VUgras{baleine}
\textsuperscript{(39)} et un \VUgras{requin} \textsuperscript{(40)} en baudruche.

%%K t16-07-1 p
7. --- Je ne me suis pas arrêté longtemps à les regarder, parce que j'ai aperçu
ce dont je rêvais : un magnifique \VUgras{théâtre de marionnettes}
\textsuperscript{(41)} avec de splendides \VUgras{décors} et un ravissant
\VUgras{rideau} rouge. Sur la \VUgras{scène}, deux acteurs semblaient jouer un
\VUgras{rôle} dans une \VUgras{pièce} très intéressante, parce que les petits
\VUgras{spectateurs} de carton installés dans les loges, ou assis dans les
\VUgras{fauteuils} et au \VUgras{parterre}, derrière le \VUgras{chef
d'orchestre}, applaudissaient avec vivacité.

%%K t16-07-2 p
Malheureusement, ce théâtre coûtait très cher.

%%K t16-08-1 p
8. --- D'ailleurs, il était difficile de choisir un jouet, parce que les
vitrines étaient remplies de jouets de toutes sortes : un \VUgras{Arlequin}
\textsuperscript{(42)}, un \VUgras{Polichinelle} \textsuperscript{(43)}, un
\VUgras{Pierrot} \textsuperscript{(44)}, des \VUgras{hiboux}
\textsuperscript{(45)} battant des ailes, des \VUgras{merles}
\textsuperscript{(46)} siffleurs, des \VUgras{caniches} aboyeurs, un
\VUgras{phonographe} \textsuperscript{(47)} et des \VUgras{jeux de patience}. Un
de ces jeux m'a beaucoup amusé : il représentait un \VUgras{combat naval}
\textsuperscript{(48)} dans lequel un \VUgras{navire} était attaqué par des
\VUgras{pirates} près d'une terre plantée de \VUgras{cocotiers}.

%%K t16-noto-1 noto
\VUnotes{143.2mm}{%
\textit{(*) Voir le tableau n\textsuperscript{o} 6.}%
}

%%K t16-09-1 p
9. --- Je pouvais contempler très à l'aise toutes ces merveilles, parce qu'il y
avait peu de monde dans le magasin : à peine quelques flâneurs, un
\VUgras{dragon} \textsuperscript{(49)} et un \VUgras{encaisseur}
\textsuperscript{(50)} qui présentait une \VUgras{traite} à la
\VUgras{caisse principale} \textsuperscript{(51)}.

%%K t16-10-1 p
10. --- J'ai demandé à mon grand-père à quoi servaient les livres posés sur des
tablettes derrière la \VUgras{caissière}; il m'a répondu que c'étaient des
\VUgras{livres de comptabilité}.

%%K t16-tit-4 sub
\VUsaut{3.6mm}
\VUpk{10.2pt}{40}{On regarde les vitrines.}
\VUsaut{2.8mm}

%%K t16-11-1 p
11. --- En quittant le rayon des jouets, nous sommes allés à la sellerie jeter
un coup d'œil sur les \VUgras{selles} \textsuperscript{(53)}, les
\VUgras{étriers} \textsuperscript{(54)}, les \VUgras{brides}
\textsuperscript{(55)}, les \VUgras{mors} \textsuperscript{(56)} et les
\VUgras{cravaches}. Pendant que le \VUgras{chef de rayon} \textsuperscript{(57)}
donnait quelques renseignements à mon grand-père, je suis allé regarder
l'étalage des \VUgras{porcelaines} et des \VUgras{cristaux}
\textsuperscript{(58)}, où il y a de beaux \VUgras{saladiers} et de délicates
\VUgras{théières} \textsuperscript{(59)}.

%%K t16-12-1 p
12. --- Pour monter au premier étage, nous sommes passés derrière la vitrine des
\VUgras{corsets} \textsuperscript{(60)}, à côté de la bonneterie, où étaient
étalés de très beaux \VUgras{caleçons} \textsuperscript{(63)}. Tout près, une
\VUgras{vendeuse} \textsuperscript{(61)} faisait choisir à une \VUgras{cliente}
\textsuperscript{(62)} de belles \VUgras{étoffes de soie} \textsuperscript{(64)}.

%%K t16-12-2 p
En montant l'escalier, j'ai remarqué que le magasin était décoré de
\VUgras{lanternes} \textsuperscript{(65)} japonaises, d'\VUgras{ombrelles}
\textsuperscript{(66)} et d'immenses \VUgras{palmiers} \textsuperscript{(70)}.

%%K t16-13-1 p
13. --- Au premier étage, il n'y avait pas grand-chose à voir : seulement des
meubles (*), des \VUgras{coffres-forts} \textsuperscript{(67)}, des
\VUgras{commodes} \textsuperscript{(68)} et des \VUgras{poussettes}
\textsuperscript{(69)}. Mais la vue d'ensemble du magasin était très
\VUgras{amusante}.

%%K t16-14-1 p
14. --- À nos pieds, j'ai vu les instruments de musique : des \VUgras{harpes}
\textsuperscript{(71)}, des \VUgras{guitares} \textsuperscript{(72)}, des
\VUgras{mandolines} \textsuperscript{(73)}, des \VUgras{cornets à pistons}
\textsuperscript{(74)}, des \VUgras{clarinettes} \textsuperscript{(75)}, des
violons avec leur \VUgras{archet} \textsuperscript{(76)} et même une
\VUgras{grosse caisse} \textsuperscript{(77)}. Un peu plus loin, au rayon de la
papeterie, un \VUgras{étudiant} \textsuperscript{(78)} marchandait une serviette
au \VUgras{commis} \textsuperscript{(79)}. Près d'eux, j'ai distingué un bel
\VUgras{abécédaire} \textsuperscript{(80)}.

%%K t16-14-2 p
Au milieu du magasin, on avait dressé un grand \VUgras{arbre de Noël}
\textsuperscript{(81)} tout garni de chandelles, de babioles et de jouets de
toutes sortes : des \VUgras{chevaux de bois} \textsuperscript{(82)}, des
\VUgras{poupées} \textsuperscript{(83)}, etc.

%%K t16-14-3 p
La \VUgras{parfumerie} \textsuperscript{(84)}, avec ses \VUgras{dentifrices} et
ses \VUgras{parfums} variés, répandait une très bonne odeur qui montait
jusqu'à nous.

%%K t16-tit-5 sub
\VUsaut{3.4mm}
\VUpk{10.2pt}{40}{Nous achetons quelque chose.}
\VUsaut{2.8mm}

%%K t16-15-1 p
15. --- Comme mon grand-père commençait à trouver le temps long, nous sommes
redescendus par le grand escalier. Il s'est arrêté un moment devant un
\VUgras{tableau d'astronomie} \textsuperscript{(85)} qui représentait, m'a-t-il
dit, des \VUgras{comètes} \textsuperscript{(a)}, des \VUgras{éclipses de lune et
de soleil} \textsuperscript{(b)}, une rose des vents avec les \VUgras{quatre
points cardinaux} \textsuperscript{(c)} \VUgras{(nord, sud, est et ouest)}, une
carte des deux \VUgras{hémisphères} \textsuperscript{(d)}, les
\VUgras{planètes} \textsuperscript{(e)} et le \VUgras{système solaire}
\textsuperscript{(f)}. Je n'ai rien compris à tout cela, et j'étais bien plus
intéressé par la livrée galonnée d'un \VUgras{garçon livreur}
\textsuperscript{(86)} qui passait, et par les beaux \VUgras{éventails}
\textsuperscript{(87)} qui étaient près de moi, à côté des
\VUgras{porte-monnaie} \textsuperscript{(88)} et des \VUgras{portefeuilles}
\textsuperscript{(89)}.

%%K t16-16-1 p
16. --- J'ai admiré aussi, sur l'étalage, des \VUgras{plumes}
\textsuperscript{(90)} et des \VUgras{boucles de ceinture}
\textsuperscript{(91)}, et les belles \VUgras{fleurs artificielles}
\textsuperscript{(94)} gracieusement disposées dans des paniers. Les
\VUgras{violettes}, les \VUgras{giroflées}, les \VUgras{jacinthes}
\textsuperscript{(95)}, les \VUgras{géraniums} \textsuperscript{(96)}, les
\VUgras{lis} \textsuperscript{(97)} et les \VUgras{capucines}
\textsuperscript{(98)} étaient très bien imitées.

%%K t16-tit-6 sub
\VUsaut{4.4mm}
\VUpk{10.2pt}{40}{Le cadeau de mon grand-père.}
\VUsaut{2.8mm}

%%K t16-17-1 p
17. --- J'ai demandé à mon grand-père de m'acheter une des \VUgras{brosses à
dents} \textsuperscript{(92)} qui étaient dans une boîte, à côté des
\VUgras{épingles à cheveux} \textsuperscript{(93)}. Il a accepté et je suis allé
moi-même payer mon achat à la caisse, près de laquelle on préparait un important
\VUgras{envoi} \textsuperscript{(100)} pour un \VUgras{client}
\textsuperscript{(99)}.

%%K t16-18-1 p
18. --- Tout à coup, il y a eu un léger désordre parmi les personnes arrêtées
près des \VUgras{bronzes d'art} \textsuperscript{(101)}. Un \VUgras{voleur}
\textsuperscript{(102)} venait d'être pris en flagrant délit par un
\VUgras{inspecteur} \textsuperscript{(103)}, pendant qu'il essayait de
dévaliser quelqu'un. On a pris soin de faire venir un policier pour
l'\VUgras{arrêter} et, au moment où nous sortions, le malfaiteur était emmené en
prison.
\VUsaut{6.0mm}
\VUfilet{22mm}
